% !TeX root = ../sections/02_exercises.tex
\subsection{2_3}
\begin{exercise}[情報数学 A 中間試験 2]
	閉区間 \([-1,1]\) 上の関数
	\[
	\psi_1(x)=x,\qquad
	\psi_2(x)=x^2,\qquad
	\psi_3(x)=1
	\]
	を考え，関数 \(\psi_k(x)\) と \(\psi_l(x)\) の内積を
	\[
	(\psi_k,\psi_l)
	=
	\int_{-1}^{1}\psi_k(x)\psi_l(x)\,dx
	\]
	で定める．
	
	正規化を伴わないシュミットの直交化法により，
	関数列
	\[
	[\psi_1,\psi_2,\psi_3]
	\]
	から直交関数列
	\[
	[\varphi_1,\varphi_2,\varphi_3]
	\]
	を作れ．
\end{exercise}
\begin{background}
	この問題では，閉区間 \([-1,1]\) 上の関数
	\[
	\psi_1(x)=x,\qquad
	\psi_2(x)=x^2,\qquad
	\psi_3(x)=1
	\]
	を考える。
	
	関数同士の内積は，
	\[
	(\psi_k,\psi_l)
	=
	\int_{-1}^{1}\psi_k(x)\psi_l(x)\,dx
	\]
	で定義されている。
	
	これは，ベクトルの内積
	\[
	\vec{u}\cdot \vec{v}
	\]
	の関数版である。
	ベクトルの場合，二つのベクトルが直交するとは内積が \(0\) になることである。
	同様に，関数 \(f,g\) が直交するとは，
	\[
	(f,g)=0
	\]
	となることをいう。
	
	この問題の目的は，関数列
	\[
	[\psi_1,\psi_2,\psi_3]
	\]
	から，互いに直交する関数列
	\[
	[\varphi_1,\varphi_2,\varphi_3]
	\]
	を作ることである。
	
	そのために，シュミットの直交化法を使う。
	
	正規化を伴わないシュミットの直交化法では，
	長さを \(1\) にする必要はない。
	したがって，求める関数列は
	\[
	(\varphi_i,\varphi_j)=0
	\quad (i\neq j)
	\]
	を満たせばよい。
	
	一方，正規化を伴う場合は，さらに
	\[
	(\varphi_i,\varphi_i)=1
	\]
	となるように，各関数をノルムで割る。
	
	この問題では「正規化を伴わない」と書かれているので，
	互いに直交する関数を作れば十分である。
	
	シュミットの直交化法の基本思想は，
	新しい関数から，すでに作った直交関数の方向成分を取り除くことである。
	
	ベクトルでいえば，
	\[
	\text{新しいベクトル}
	-
	\text{前のベクトル方向への射影}
	\]
	を行っている。
	関数の場合も同じで，内積を用いて射影成分を計算する。
\end{background}
\begin{strategy}
	正規化を伴わないシュミットの直交化法では，次の順番で進める。
	
	まず，
	\[
	\varphi_1=\psi_1
	\]
	とおく。
	
	この問題では，
	\[
	\psi_1(x)=x
	\]
	なので，
	\[
	\varphi_1(x)=x
	\]
	である。
	
	次に，\(\psi_2\) から \(\varphi_1\) 方向の成分を引く。
	すなわち，
	\[
	\varphi_2
	=
	\psi_2
	-
	\frac{(\psi_2,\varphi_1)}{(\varphi_1,\varphi_1)}
	\varphi_1
	\]
	とする。
	
	この問題では，
	\[
	(\psi_2,\varphi_1)
	=
	(x^2,x)
	=
	\int_{-1}^{1}x^3\,dx
	\]
	である。
	\(x^3\) は奇関数なので，
	\[
	\int_{-1}^{1}x^3\,dx=0
	\]
	である。
	
	したがって，
	\[
	\varphi_2=x^2
	\]
	となる。
	
	最後に，\(\psi_3\) から，\(\varphi_1\) 方向の成分と
	\(\varphi_2\) 方向の成分を引く。
	すなわち，
	\[
	\varphi_3
	=
	\psi_3
	-
	\frac{(\psi_3,\varphi_1)}{(\varphi_1,\varphi_1)}
	\varphi_1
	-
	\frac{(\psi_3,\varphi_2)}{(\varphi_2,\varphi_2)}
	\varphi_2
	\]
	である。
	
	この問題では，
	\[
	\psi_3=1,\qquad
	\varphi_1=x,\qquad
	\varphi_2=x^2
	\]
	なので，
	\[
	\varphi_3
	=
	1
	-
	\frac{(1,x)}{(x,x)}x
	-
	\frac{(1,x^2)}{(x^2,x^2)}x^2
	\]
	となる。
	
	まず，
	\[
	(1,x)=\int_{-1}^{1}x\,dx=0
	\]
	である。
	
	次に，
	\[
	(1,x^2)=\int_{-1}^{1}x^2\,dx=\frac{2}{3}
	\]
	である。
	
	さらに，
	\[
	(x^2,x^2)=\int_{-1}^{1}x^4\,dx=\frac{2}{5}
	\]
	である。
	
	したがって，
	\[
	\frac{(1,x^2)}{(x^2,x^2)}
	=
	\frac{2/3}{2/5}
	=
	\frac{5}{3}
	\]
	となる。
	
	よって，
	\[
	\varphi_3
	=
	1-\frac{5}{3}x^2
	\]
	である。
	
	したがって，正規化を伴わない直交関数列として，
	\[
	\varphi_1=x,\qquad
	\varphi_2=x^2,\qquad
	\varphi_3=1-\frac{5}{3}x^2
	\]
	を得る。
	
	また，直交性は定数倍しても変わらないので，
	\[
	\varphi_3=3-5x^2
	\]
	としてもよい。
	その場合，
	\[
	\varphi_1=x,\qquad
	\varphi_2=x^2,\qquad
	\varphi_3=3-5x^2
	\]
	と書ける。
	
	この問題の発想は，
	\[
	\text{内積で射影成分を測る}
	\Rightarrow
	\text{前に作った関数方向の成分を引く}
	\Rightarrow
	\text{直交関数列を作る}
	\]
	という流れである。
\end{strategy}